% GENERATED FILE. Edit canonical lessons, then run npm run generate:books.
% canonical-source-hash: fnv1a64:77aae4f0cbac01f6
% canonical-lessons: TE-C308-gunta, TE-C308-mukka, TE-C308-ancu, TE-C308-mudi, TE-C308-varusa

\chapter{Pit, Piece, Edge, Knot, Row}
\label{ch:te-pit-piece-edge-knot-row}
\hlchaptermodality{308}

\begin{chapteropening}
\textbf{By the end of this chapter:} \emph{I can say a pit, a piece, an edge, a knot and a row.}

Complete the last lesson of chapter 308: I can say a pit, a piece, an edge, a knot and a row.
\end{chapteropening}

\section[\texorpdfstring{\te{గుంట} (guṇṭa)}{గుంట (guṇṭa)}]{\te{గుంట} (guṇṭa) — a pit, a pothole}

\label{lesson:TE-C308-gunta}



\begin{quote}
Before the new one: say the Telugu for an incense stick, then the Telugu for kumkum.
\end{quote}

\subsection*{You'll want to know: \te{గుంట}}
\textbf{\te{గుంట}} — \emph{guṇṭa} — \textquotedblleft{}a pit, a pothole\textquotedblright{}.

\begin{grammarlens}[title={What you've built}]
One more word for this chapter.
\end{grammarlens}

\subsection*{Your turn}
\begin{itemize}
\raggedright
  \item \emph{Say it:} \emph{guṇṭa}
  \item \emph{Say it:} \emph{guṇṭa}, once more
  \item \emph{Say it:} say \emph{guṇṭa}
  \item \emph{Recall:} say the Telugu for an incense stick, then the Telugu for kumkum, then say \emph{guṇṭa} again
\end{itemize}

\begin{tcolorbox}[breakable,colback=teal!4,colframe=teal!35!black,title={Before you move on}]
What does \te{గుంట} mean? (\textquotedblleft{}A pit, a pothole\textquotedblright{}.) Say it once more.
\end{tcolorbox}
\section[\texorpdfstring{\te{ముక్క} (mukka)}{ముక్క (mukka)}]{\te{ముక్క} (mukka) — a piece}

\label{lesson:TE-C308-mukka}



\begin{quote}
Before the new one: say the Telugu for kumkum, then the Telugu for a pit.
\end{quote}

\subsection*{You'll want to know: \te{ముక్క}}
\textbf{\te{ముక్క}} — \emph{mukka} — \textquotedblleft{}a piece\textquotedblright{}.

\begin{grammarlens}[title={What you've built}]
One more word for this chapter.
\end{grammarlens}

\subsection*{Your turn}
\begin{itemize}
\raggedright
  \item \emph{Say it:} \emph{mukka}
  \item \emph{Say it:} \emph{mukka}, once more
  \item \emph{Say it:} say \emph{mukka}
  \item \emph{Recall:} say the Telugu for kumkum, then the Telugu for a pit, then say \emph{mukka} again
\end{itemize}

\begin{tcolorbox}[breakable,colback=teal!4,colframe=teal!35!black,title={Before you move on}]
What does \te{ముక్క} mean? (\textquotedblleft{}A piece\textquotedblright{}.) Say it once more.
\end{tcolorbox}
\section[\texorpdfstring{\te{అంచు} (añcu)}{అంచు (añcu)}]{\te{అంచు} (añcu) — an edge, a rim}

\label{lesson:TE-C308-ancu}



\begin{quote}
Before the new one: say the Telugu for a pit, then the Telugu for a piece.
\end{quote}

\subsection*{You'll want to know: \te{అంచు}}
\textbf{\te{అంచు}} — \emph{añcu} — \textquotedblleft{}an edge, a rim\textquotedblright{}.

\begin{grammarlens}[title={What you've built}]
One more word for this chapter.
\end{grammarlens}

\subsection*{Your turn}
\begin{itemize}
\raggedright
  \item \emph{Say it:} \emph{añcu}
  \item \emph{Say it:} \emph{añcu}, once more
  \item \emph{Say it:} say \emph{añcu}
  \item \emph{Recall:} say the Telugu for a pit, then the Telugu for a piece, then say \emph{añcu} again
\end{itemize}

\begin{tcolorbox}[breakable,colback=teal!4,colframe=teal!35!black,title={Before you move on}]
What does \te{అంచు} mean? (\textquotedblleft{}An edge, a rim\textquotedblright{}.) Say it once more.
\end{tcolorbox}
\section[\texorpdfstring{\te{ముడి} (muḍi)}{ముడి (muḍi)}]{\te{ముడి} (muḍi) — a knot}

\label{lesson:TE-C308-mudi}



\begin{quote}
Before the new one: say the Telugu for a piece, then the Telugu for an edge.
\end{quote}

\subsection*{You'll want to know: \te{ముడి}}
\textbf{\te{ముడి}} — \emph{muḍi} — \textquotedblleft{}a knot\textquotedblright{}.

\begin{grammarlens}[title={What you've built}]
One more word for this chapter.
\end{grammarlens}

\subsection*{Your turn}
\begin{itemize}
\raggedright
  \item \emph{Say it:} \emph{muḍi}
  \item \emph{Say it:} \emph{muḍi}, once more
  \item \emph{Say it:} say \emph{muḍi}
  \item \emph{Recall:} say the Telugu for a piece, then the Telugu for an edge, then say \emph{muḍi} again
\end{itemize}

\begin{tcolorbox}[breakable,colback=teal!4,colframe=teal!35!black,title={Before you move on}]
What does \te{ముడి} mean? (\textquotedblleft{}A knot\textquotedblright{}.) Say it once more.
\end{tcolorbox}
\section[\texorpdfstring{\te{వరుస} (varusa)}{వరుస (varusa)}]{\te{వరుస} (varusa) — a row, a queue}

\label{lesson:TE-C308-varusa}



\begin{quote}
Before the new one: say the Telugu for an edge, then the Telugu for a knot.
\end{quote}

\subsection*{You'll want to know: \te{వరుస}}
\textbf{\te{వరుస}} — \emph{varusa} — \textquotedblleft{}a row, a queue\textquotedblright{}.

\begin{grammarlens}[title={What you've built}]
One more word for this chapter.
\end{grammarlens}

\subsection*{Your turn}
\begin{itemize}
\raggedright
  \item \emph{Say it:} \emph{varusa}
  \item \emph{Say it:} \emph{varusa}, once more
  \item \emph{Say it:} say \emph{varusa}
  \item \emph{Recall:} say the Telugu for an edge, then the Telugu for a knot, then say \emph{varusa} again
\end{itemize}

\begin{tcolorbox}[breakable,colback=teal!4,colframe=teal!35!black,title={Before you move on}]
What does \te{వరుస} mean? (\textquotedblleft{}A row, a queue\textquotedblright{}.) Say it once more.
\end{tcolorbox}
